\scanpage{8}\mspage{6}
avec une relation d'associativité, disant que
\diagram{\begin{tikzcd}[ampersand replacement=\&]
P_0*P_{\unsure{1|3}}*P_2 \arrow[r] \arrow[d] \& Q_0*P_2 \arrow[d]\\
P_0*Q_1 \arrow[r] \& R
\end{tikzcd}}
est commutatif, alors on en déduit \deleted{trois} \inserted{deux} relations de transitivité et une de commutation, pour les opérations associées aux éléments de $\pi_0(P_0),\pi_0(P_1),\pi_0(P_2)$, qu'on \unsure{n'explicitera} pas, sauf l'une :
\[
\text{Soit }\beta\in\deleted{\unsure{\pi_0}}P_1,\ \gamma\unsure{\in|:}P_2
\]
\[
\begin{aligned}
\text{d'où par }\beta:\quad &c_{u_1}(\beta):P_0\longrightarrow Q_0\\
\text{et par }\gamma:\quad &c_{u_4}(\gamma):Q_0\longrightarrow R\\
\text{et par }\underset{\substack{\rotatebox{-90}{$\in$}\\\pi_0(Q_1)}}{\deleted{\unsure{\beta\gamma}}\inserted{\beta\underset{u_2}{*}\gamma}}:\quad &c_{u_3}(\beta\gamma):P_0\longrightarrow\deleted{\unsure{Q}}R
\end{aligned}
\]
\unsure{eh|et} bien on a
\[
\boxed{c(\beta\underset{u_2}{*}\gamma)=c(\gamma)c(\beta)}.
\]
Il faut encore définir $\beta\underset{u_2}{*}\gamma$, via l'accouplement $P_1*P_2\xrightarrow{u_2}Q_1$, on a en effet
\[
\begin{aligned}
\pi_0(P_1)\times\pi_0(P_{\unsure{7|2|1}})&\longrightarrow\pi_0(P_1*P_2)\\
(\beta,\gamma)&\longmapsto\beta*\gamma
\end{aligned}
\]
et on compose avec
\[
\pi_0(u_2):P_1*P_2\longrightarrow Q_1\qquad\unsure{?|:}
\]
% archive sheet number: 7

